\section{Feasibility, Risk and Engineering Budgets}\label{mission-completion-requirements}
\subsection{Mission Feasibility}\label{sec:mission-feasibility}
The engineering studies support ANTARA's concept-level operational and resource feasibility for the evaluated scenarios. The corrected integrated simulation maintains positive recorded battery energy while scheduling startup, attitude transitions, science acquisition, ground relay and recovery over a complete year. Its time-aligned data-service calculation accommodates the sampled science-record workload. The mission-level finding is \textbf{CONDITIONALLY SUPPORTED}: detector response, selected-actuator capability, component thermal limits and complete mass allocation still determine whether these modeled opportunities can deliver the approved science objectives. Hardware qualification is not a prerequisite for the scenario-level findings established here.

\paragraph{Scientific feasibility --- CONDITIONALLY SUPPORTED.}
The 5--100~MeV objective, spatial mapping tiers and uncertainty criteria remain unchanged. Radiation maps and incident-rate calculations provide a quantitative sizing basis, and the 292-byte histogram record supports compact data return. Ideal geometric factors and unit-efficiency incident rates do not establish accepted detector counts or calibrated flux. Energy-dependent response, background discrimination, effective geometric factor and reconstruction uncertainty remain essential concept questions. MMSC, MSC and EMSC achievement is not yet assessed from accepted science observations.

\paragraph{Orbital feasibility --- CONDITIONALLY SUPPORTED.}
The approximately 600~km, $97.6^\circ$ orbit studies provide SAA acquisition and ground-contact opportunities. Seven RAAN cases quantify sensitivity to launch-dependent illumination, with average generation spanning approximately 9.02--11.17~W. These results support the orbit concept under the respective study assumptions; insertion RAAN remains uncontrolled, and the studies do not represent a single common 2029 configuration. Passive decay is assessed separately below.

\paragraph{EPS feasibility --- CONDITIONALLY SUPPORTED.}
The corrected annual scenario records 50,288.809~Wh generation, 30,866.103~Wh modeled demand and a minimum battery energy of 15.098~Wh. These results support sustained energy supply for the evaluated schedule. Application to the planned 25~Wh battery requires selected-cell limits, startup reserve, peak electrical demand and consistent array/EOL assumptions; the proposed 28.8~Wh rated pack is not substituted into the calculation.

\paragraph{ADCS feasibility --- CONDITIONALLY SUPPORTED.}
The integrated case completes a 2,460~s detumble event and allocates time to Sun Pointing, Detection, Relay and slews. This supports operational sequencing. The separate synthetic Monte Carlo result of 16,133 successes in 50,000 trials is incompatible with the current initial-condition model and does not establish reliability of the intended reaction-wheel/magnetorquer system. Pointing accuracy, settling time and torque/momentum authority remain unassessed for the selected architecture.

\paragraph{Thermal feasibility --- NOT YET ASSESSED against component limits.}
The seven-node model predicts internal temperatures of $-14.061$ to $21.159^\circ$C and external-face temperatures of $-25.958$ to $31.239^\circ$C for its independent heat schedule. This establishes a preliminary thermal envelope. Component compliance cannot be inferred without operating/survival limits and coupled hot/cold cases; no thermal margin is assigned.

\paragraph{Communications feasibility --- CONDITIONALLY SUPPORTED.}
The annual study provides 820 contacts and 178,682~s of visibility intersecting telemetry-labelled Relay. At the full ideal 19.2~kbps gross rate, calculated capacity is 428,836,800 bytes, accommodating the 46,637,364-byte sampled workload with zero final backlog. This supports scheduled data return; the RF link budget and actual net science throughput remain unassessed for the selected implementation.

\paragraph{OBC and data-handling feasibility --- CONDITIONALLY SUPPORTED.}
The reference codec and queue calculations support compact histogram assembly and storage. The 75,336-byte nominal queue fits the candidate 256~KiB science RAM allocation before metadata. The seven-day outage peak of 1,010,028 bytes is \textbf{NOT SUPPORTED BY CURRENT RESULTS} for RAM-only retention; persistent spill is required. The candidate 128~MiB allocation accommodates this modeled outage before reserves. Total OBC memory, execution demand and the accepted retention envelope remain separate allocations.

\paragraph{Mission operations feasibility --- SUPPORTED BY SIMULATION for the evaluated schedule.}
The corrected LEOP 24-hour, No-LEOP 24-hour and LEOP one-year scenarios demonstrate startup-to-operations sequencing, repeated acquisition and ground-relay opportunities, and sustained modeled energy availability. The annual record includes 332,882~s in Safe Mode, so the result includes interruptions rather than assuming continuous science. Actual station availability, accepted detector observations and unresolved transition-duration semantics constrain application beyond the evaluated scenario. The source epoch and assumptions are stated once in Section~\ref{sec:sim-leop}.

\paragraph{Mechanical and mass feasibility --- NOT YET ASSESSED for the complete spacecraft.}
The 3U arrangement and defined FPDU geometry establish a packaging concept. The 1536.3~g specified-mass subtotal excludes essential allocations and cannot establish compliance with the 3~kg design limit. Complete payload/deployment mass, panel quantity, packaged clearances and launch-provider loads remain necessary to assess mechanical feasibility.

\subsection{Measures of Performance and Effectiveness}\label{measures-of-performance-mops}
Measures of Performance cover battery state, Detection pointing, science-data return, payload livetime and mission lifetime. Measures of Effectiveness cover proton measurements, data return, radiation mapping and model comparison. Appendix~\ref{app:source-05} retains the original thresholds and goals; Section~\ref{mission-success-criteria} remains authoritative for science success. Operational opportunity and queue capacity support these objectives without replacing their calibrated-measurement criteria.

\subsection{Engineering Budgets}\label{sec:budget-closure}
The specified component masses total 1536.3~g against a design target below 3~kg. Payload, deployment hardware and panel-assembly quantity remain incomplete, so no mass margin is reported. Component costs remain TBD. Energy planning uses 25~Wh; the proposed Type A pack is rated 28.8~Wh. Science handling uses the 292-byte reference record and the corrected time-aligned workload, with RAM and persistent storage assessed separately above. Detailed component and budget calculations are retained in Appendices~\ref{app:source-05} and~\ref{app:detailed-budgets}.

\subsection{Operational and Engineering Risks}\label{risk-method-and-register}
Deployment failure, unsuccessful detumbling, loss of pointing, radiation effects, communications outages, battery degradation and thermal excursion remain mission risks. The simulations quantify selected consequences: Safe Mode reduces acquisition opportunity, a seven-day communications outage exceeds the candidate RAM queue, and launch RAAN changes generation. The available risk register does not provide complete likelihood/consequence scores; these effects are retained without assigning unsupported reliability or margin. The detailed register remains in Appendix~\ref{app:source-05}.

\subsection{End-of-Life and Debris Mitigation}\label{end-of-life-deorbit-and-debris-mitigation-compliance}
ANTARA has no propulsion and relies on passive orbital decay. The approximately five-year operational science objective is distinct from the modeled 8.2--10.0-year drag-study endpoints across seven RAAN cases and two runs per case. Most traces end near 300~km; one ends near 335~km. These are propagation cutoffs, not atmospheric reentry. Disposal feasibility requires continuation to a physical disposal endpoint with applicable epoch and aerodynamic assumptions. Appendix~\ref{app:updated-raan} retains the study details.

\subsection{Remaining Concept-Level Questions}\label{sec:open-issues}
The remaining questions affecting mission feasibility are detector acceptance and reconstruction over the required energy range; mixed-species trigger demand and end-to-end timing; compatible ADCS authority and pointing allocations; selected-cell limits and peak/reserve power allocations; component thermal limits and coupled operating cases; complete mass and mechanical allocations; and an accepted RF service/outage envelope. Science coverage also requires eligible-cell, independent-pass and uncertainty definitions, including resolution of the existing MSC channel-count inconsistency. Fragmented acquisition can exceed the preferred three contiguous readings per visit and requires an explicit record policy. A common configuration and consistent transition timing are needed before combining separate studies into mission-wide margins. These are substantive constraints, not routine qualification prerequisites.

\subsection{Further Design Development and PDR Scope}\label{sec:mcr-pdr-scope}
PDR development will refine the RF link budget and net throughput including protocol overhead and retransmission; finalize detector materials, optical response, effective geometric factor and reconstruction; complete mass, structural, deployment and packaged-interface design; define component thermal/environmental qualification; and implement OBC timing, memory, persistent storage, command and subsystem interfaces. These activities build on the demonstrated scenario-level feasibility. The concept questions above remain explicit until resolved by suitable engineering evidence and are not deferred merely as qualification work.
